\chaptertopic{Process Architecture}

\begin{learningobjectives}
  \item Define process architecture.
  \item Distinguish primary, management, and supporting processes.
  \item Explain how a process architecture supports analysis.
\end{learningobjectives}

\section{Organizing Processes}

A \term{process architecture} is an organized view of the processes in an
enterprise and the relationships between them. It helps people see how work
across teams contributes to organizational goals.

\definition[def:process-architecture]{Process Architecture}{
A structured overview of an organization's processes and their relationships.
It may group processes into primary, management, and supporting categories.
}

\begin{center}
  \begin{tabularx}{\textwidth}{@{}lYY@{}}
    \toprule
    Process type & Main purpose & Example \\
    \midrule
    Primary & Deliver value to customers & Fulfil a customer order \\
    Management & Direct and monitor the organization & Set business goals \\
    Supporting & Enable other processes & Maintain information systems \\
    \bottomrule
  \end{tabularx}
\end{center}

\keyconcept{End-to-end perspective}{
A process architecture makes dependencies visible across departments. Analyze
how the complete process delivers an outcome rather than optimizing each team
in isolation.
}

\section{Representing a Workflow}

Programming and technical subjects can use the same template. The code block
below illustrates an order validation step.

\begin{codeblock}{C\#}
public class Order
{
    public int Id { get; set; }

    public bool IsValid()
    {
        return Id > 0;
    }
}
\end{codeblock}

Other supported code languages include JavaScript, TypeScript, Python, SQL,
HTML, CSS, and Bash. Use the spelling shown in the language name.

\section{Revision Questions}

\begin{questionlist}
  \q{What is process architecture?}

  \qa[q:primary-process]{What distinguishes a primary process from a
  supporting process?}
  {A primary process delivers value directly to customers; a supporting
  process supplies capabilities or resources that enable other processes.}

  \examqa{Explain how a process architecture can reveal opportunities for
  improvement.}
  {It makes relationships and handoffs visible, helping identify duplication,
  bottlenecks, unclear responsibilities, and gaps between processes.}
\end{questionlist}

See Definition~\ref{def:process-architecture} and Question~
\ref{q:primary-process} when reviewing the process categories.

\important{
A process category describes the role a process plays in the organization;
it does not indicate that one category is inherently more important than
another.
}

\begin{summarybox}
\begin{itemize}
  \item Process architecture organizes processes and shows their connections.
  \item Primary processes deliver customer value.
  \item Management and supporting processes enable effective delivery.
\end{itemize}
\end{summarybox}
